% Created 2026-04-26 dom 18:26
% Intended LaTeX compiler: xelatex
\documentclass[11pt]{article}
\usepackage{graphicx}
\usepackage{longtable}
\usepackage{wrapfig}
\usepackage{rotating}
\usepackage[normalem]{ulem}
\usepackage{capt-of}
\usepackage{hyperref}
\author{Nícolas Rosenthal Dal Corso}
\date{\date}
\title{anotações para a resolução do laboratório}
\hypersetup{
 pdfauthor={Nícolas Rosenthal Dal Corso},
 pdftitle={anotações para a resolução do laboratório},
 pdfkeywords={},
 pdfsubject={},
 pdfcreator={Emacs 30.2 (Org mode 9.7.11)}, 
 pdflang={English}}
\begin{document}

\maketitle
\tableofcontents

\section{Dificuldades e Erros}
\label{sec:orga6bcc0a}
{[}modelo /baseline/] (\texttt{n-rosenthal-{}-{}submission-1})
    ->
diversos modelos para testar e aprender técnicas avançadas
    ->
{[}modelo /stacking-ensemble 4/] (\texttt{n-rosenthal-{}-{}submission-3})


No modelo \emph{baseline}:
\begin{enumerate}
\item Após definir um modelo \emph{baseline}, percebi que tratei a coluna \texttt{Id} como outra coluna qualquer;
\item Na primeira vez, eu defini duas estratégias de imputação de valores ausentes, uma depois da outra. Isto ocasionou uma \emph{dupla imputação}. Isso não faz sentido nenhum.
\end{enumerate}
\section{Modelo \emph{Baseline}}
\label{sec:org883ef71}
\subsection{Pré-Processamento dos Dados: Limpeza}
\label{sec:org5404ad6}
\subsubsection{Importação das Bibliotecas}
\label{sec:org61bba53}

\begin{verbatim}
import numpy as np
import pandas as pd
\end{verbatim}
\subsubsection{Leitura dos Dados}
\label{sec:orgf34f21a}
Ler os dados a partir de documento \texttt{.csv} e inspeção de sua dimensionalidade. Também vamos armazenar o numéro de colunas em uma variável de trabalho.

\begin{verbatim}
train_df: pd.DataFrame = pd.read_csv(r"./lab05_files/lab05_files/train_student.csv");
test_df: pd.DataFrame  = pd.read_csv(r"./lab05_files/lab05_files/test_student.csv");

print("Dimensões:");
print("\tdados de treinamento: ", train_df.shape);
print("\tdados de teste      : ", test_df.shape);

n_cols: int = test_df.shape[1];
\end{verbatim}

Nosso objetivo é, a partir das \texttt{81}  colunas, predizer o valor da coluna \texttt{SalesPrice}. É observável também que os dados de teste possuem uma coluna a mais do que os dados de treinamento.
\subsubsection{Diferença na Quantidade de Colunas}
\label{sec:org94c0d02}
\begin{verbatim}
train_cols = set(train_df.columns)
test_cols  = set(test_df.columns)

missing_in_test = train_cols - test_cols
extra_in_test   = test_cols - train_cols

print("Colunas em `train`, mas não em `test` :", end="")
print(missing_in_test)
print("\nColunas em `test`, mas não em `train` :", end="")
print(extra_in_test)
\end{verbatim}
\begin{enumerate}
\item Natureza das Colunas
\label{sec:org6a8d5bb}
Por inspeção visual dos dados, observamos que existem colunas \emph{numéricas} e colunas \emph{categóricas}. Primeiro, vamos isolar a coluna que estamos interessados. Em seguida, vamos separar as colunas numéricas das colunas categóricas. É necessário verificar se estamos sendo capazes de extrair todas as colunas e se estamos seprando-as adequadamente.

\begin{verbatim}
feature_cols = [col for col in test_df.columns if col != 'SalesPrice']  # test_df has no SalesPrice anyway

numerical_cols = test_df[feature_cols].select_dtypes(include=['int64', 'float64']).columns
categorical_cols = test_df[feature_cols].select_dtypes(include=['object', 'category']).columns

n_cols_numerical = len(numerical_cols)
n_cols_categorical = len(categorical_cols)

print("colunas numéricas  : ", n_cols_numerical)
print("colunas categóricas: ", n_cols_categorical)
total = n_cols_numerical + n_cols_categorical
print("total              : ", total, end="")
print(" -> ", total == n_cols)   # n_cols = test_df.shape[1] = 80
\end{verbatim}

Vamos primeiro analisar questões de limpeza dos dados para tanto colunas numéricas quanto categóricas. Em seguida, analisaremos os dois conjuntos separadamente e, quando formos propriamente treinar o modelo, reuniremos as colunas.

É muito provável que existam \emph{valores ausentes} para algumas entradas. Vamos investigar:

\begin{verbatim}
# treinamento
missing_train = train_df.isnull().sum()
missing_train = missing_train[missing_train > 0].sort_values(ascending=False)

print("Valores ausentes nos dados de treinamento: ")
print(missing_train)
print(f"\ntotal: {missing_train.sum()}")

# teste
missing_test = test_df.isnull().sum()
missing_test = missing_test[missing_test > 0].sort_values(ascending=False)

print("\nValores ausentes nos dados de teste:")
print(missing_test)
print(f"\ntotal: {missing_test.sum()}")
\end{verbatim}

Muitos valores faltando, tanto em colunas categóricas quanto numéricas, e tanto para os dados teste quanto para os dados de treinamento. A resolução deverá ser aplicada para ambos conjuntos (\texttt{train, test}) de forma análoga.

Existem algumas estratégias possíveis para lidar com este problema. Acaso uma coluna tenha mais do que cinquenta por cento de valores ausentes, será mesmo que ela é importante para a predição final?

\begin{verbatim}
# Fazer uma cópia dos DataFrames para evitar perder valores
cleaned_train_df: pd.DataFrame = train_df.copy();
cleaned_test_df:  pd.DataFrame = test_df.copy();

# Definir threshold para eliminar colunas (>= 50%)
thresholds: list[float] = list(np.linspace(0.05, 0.99, num=20));
print(thresholds);

for threshold in thresholds:
    kill_cols_train = cleaned_train_df.columns[cleaned_train_df.isnull().mean() > threshold];
    kill_cols_test  = cleaned_test_df.columns[cleaned_test_df.isnull().mean() > threshold];
    kill_cols: set = set(kill_cols_test) | set(kill_cols_train)

    print("threshold `", threshold, "`\n\t ", list(kill_cols));

\end{verbatim}

Definimos um intervalo linear de \(20\) valores entre \texttt{0.05} e \texttt{0.99} para os \emph{thresholds} de forma a mostrar quais colunas possuem valores faltando. Percebems que diveras colunas possuem pelo menos 5\% dos valores ausentes: \texttt{['MasVnrType', 'GarageYrBlt', 'GarageFinish', 'PoolQC', 'Fence', 'GarageQual', 'Alley', 'MiscFeature', 'LotFrontage', 'FireplaceQu', 'GarageType', 'GarageCond']}; este é um intervalo baixo e não podemos decidir nossa análise sobre ele. O outro limite da nossa análise (\texttt{99\%} dos valores ausentes) apresenta apenas a coluna \texttt{PoolQC}.

Certamente não podemos esperar que uma coluna com valores ausentes para 99\% das entradas seja importante para predição. Alguns thresholds interssantes são \texttt{0.25, 0.50, 0.75}:

\begin{verbatim}
for threshold in [0.1, 0.25, 0.50, 0.75, 0.9]:
    kill_cols_train = cleaned_train_df.columns[cleaned_train_df.isnull().mean() > threshold];
    kill_cols_test  = cleaned_test_df.columns[cleaned_test_df.isnull().mean() > threshold];
    kill_cols: set = set(kill_cols_test) | set(kill_cols_train)

    print("threshold `", threshold, "`\n", "\t".join(list(kill_cols)));

\end{verbatim}

A decisão mais informada que tomamos é (1.) eliminar colunas com \emph{muitos} valores ausentes e (2.) decidir manualmente a partir de inspeção mais profunda, para colunas com quantidade intermediária de valores faltando.
\begin{enumerate}
\item Eliminar Colunas com Muitos Valores Ausentes
\label{sec:org1908837}
\begin{enumerate}
\item Decisão por coluna
\label{sec:orged8d3cf}
\begin{verbatim}
missing_ratio_train = train_df.isnull().mean().sort_values(ascending=False)
missing_ratio_test  = test_df.isnull().mean()

important_cols = ['MasVnrType', 'FireplaceQu', 'LotFrontage', 'GarageType']

print("Razões de ausência (train) — top 10:")
print(missing_ratio_train.head(10))
print("\nColunas que parecem ser importantes:")
for col in important_cols:
    if col in missing_ratio_train.index:
        print(f"{col}: train {missing_ratio_train[col]:.1%}, test {missing_ratio_test[col]:.1%}")
\end{verbatim}

\begin{enumerate}
\item As colunas \texttt{PoolQC}, \texttt{MiscFeature} e \texttt{Alley} são altamente esparsas (\(\geq 0.93\)). Estas podem ser removidas sem detrimento à análise.
\item \texttt{Fence} é nosso \emph{threshold} verdadeiro (\texttt{0.8}). A partir daqui, queremos que os valores ausentes sejam substituídos por uma categoria real (neste caso, ou \texttt{NoFence} ou \texttt{None} simplesmente).
\item As colunas posteriores permancem, neste momento.
\end{enumerate}
\item Limpeza Propriamente Dita
\label{sec:orga8aeac2}
\begin{verbatim}
cleaned_train = train_df.copy()
cleaned_test = test_df.copy()

# Colunas com muitos dados ausentes são eliminadas
cols_to_drop = ['PoolQC', 'MiscFeature', 'Alley', 'Fence']
cleaned_train = cleaned_train.drop(columns=cols_to_drop)
cleaned_test = cleaned_test.drop(columns=cols_to_drop)
\end{verbatim}
\end{enumerate}
\item Preencher Colunas com Valores Ausentes
\label{sec:orgb862646}
Ainda existem algumas colunas com valores ausentes e precisamos que estes sejam preenchidos para podermos avançar para as próximas etapas. Abaixo, vamos analisar coluna-a-coluna
\begin{verbatim}
# 1. Imputar valores ausentes em categorias com novas categorias
cleaned_train['MasVnrType'] = cleaned_train['MasVnrType'].fillna('None')
cleaned_test['MasVnrType'] = cleaned_test['MasVnrType'].fillna('None')

cleaned_train['FireplaceQu'] = cleaned_train['FireplaceQu'].fillna('NoFireplace')
cleaned_test['FireplaceQu'] = cleaned_test['FireplaceQu'].fillna('NoFireplace')

# *Garagem*: se ausente, então o imóvel não tem gargem
garage_cols = ['GarageType', 'GarageFinish', 'GarageQual', 'GarageCond']
for col in garage_cols:
    cleaned_train[col] = cleaned_train[col].fillna('NoGarage')
    cleaned_test[col] = cleaned_test[col].fillna('NoGarage')

# *Subsolo/Porão*: se ausente, não há subsolo
basement_cols = ['BsmtQual', 'BsmtCond', 'BsmtExposure', 'BsmtFinType1', 'BsmtFinType2']
for col in basement_cols:
    cleaned_train[col] = cleaned_train[col].fillna('NoBasement')
    cleaned_test[col] = cleaned_test[col].fillna('NoBasement')

# 2. Imputação para valores numéricos: 0, mediana, moda
# LotFrontage: mediana, a fim de evitar outliers
#              alguns imóveis podem ter `Frontage` muito grande, ou nula
median_lot = cleaned_train['LotFrontage'].median()
cleaned_train['LotFrontage'] = cleaned_train['LotFrontage'].fillna(median_lot)
cleaned_test['LotFrontage'] = cleaned_test['LotFrontage'].fillna(median_lot)

# MasVnrArea: 0 (-> sem alvenaria)
cleaned_train['MasVnrArea'] = cleaned_train['MasVnrArea'].fillna(0)
cleaned_test['MasVnrArea'] = cleaned_test['MasVnrArea'].fillna(0)

# GarageYrBlt: 0 (-> sem garagem)
cleaned_train['GarageYrBlt'] = cleaned_train['GarageYrBlt'].fillna(0)
cleaned_test['GarageYrBlt'] = cleaned_test['GarageYrBlt'].fillna(0)

# Electrical: difícil entender o que isto significa
#             -> moda, valor mais comum
mode_elec = cleaned_train['Electrical'].mode()[0]
cleaned_train['Electrical'] = cleaned_train['Electrical'].fillna(mode_elec)
cleaned_test['Electrical'] = cleaned_test['Electrical'].fillna(mode_elec)

# 3. Inspeção: ainda há valores faltando?
print(f"Valores ausentes: train = {cleaned_train.isnull().sum().sum()}, test = {cleaned_test.isnull().sum().sum()}")
if cleaned_train.isnull().sum().sum() == 0:
    print("\tLimpeza completa, não há mais valores ausentes.")
\end{verbatim}
\end{enumerate}
\end{enumerate}
\subsubsection{Pipeline para Pré-Processamento}
\label{sec:orgaf8528c}
Agora que a análise manual foi feita, podemos \textbf{refatorar} o código acima para usar estratégias mais modernas com as \emph{pipelines} do \texttt{sklearn}. Isto é interessante porque
\begin{enumerate}
\item é mais robusto e generalista;
\item elimina \emph{warnings} do tipo \texttt{FutureWarning}; e
\item faz uso das técnicas vigentes em Ciência de Dados.
\end{enumerate}

\begin{verbatim}
import pandas as pd
import numpy as np
from sklearn.impute import SimpleImputer
from sklearn.compose import ColumnTransformer

# Carregar os dados
train_df = pd.read_csv("./lab05_files/lab05_files/train_student.csv")
test_df  = pd.read_csv("./lab05_files/lab05_files/test_student.csv")


print("1. Dados carregados. Dimensões:", train_df.shape, test_df.shape)

# Recomeçamos com cópias dos datasets iniciais
cleaned_train = train_df.copy()
cleaned_test = test_df.copy()

# 1. Eliminaremos colunas com >= 0.8 valores ausentes
cols_to_drop = ['PoolQC', 'MiscFeature', 'Alley', 'Fence']
cleaned_train.drop(columns=cols_to_drop, inplace=True)
cleaned_test.drop(columns=cols_to_drop, inplace=True)

print("2. Eliminadas colunas com >= 0.8 valores ausentes.")


# 2. Separar features (X) e target (y) para treinamento
X_train = cleaned_train.drop('SalePrice', axis=1)
y_train = cleaned_train['SalePrice']
X_test = cleaned_test

print(f"3. X_train dim.: {X_train.shape}")

# 3. Converter MSSubClass para string (categórica) antes da imputação
if 'MSSubClass' in X_train.columns:
    X_train['MSSubClass'] = X_train['MSSubClass'].astype(str)
    X_test['MSSubClass'] = X_test['MSSubClass'].astype(str)
    print("4. coluna `MSSubClass` convertida para `str`")

# 4. Definir colunas para imputação com zero (feature ausente = 0)
zero_impute_cols = [
    'MasVnrArea', 'GarageYrBlt', 'GarageCars', 'GarageArea',
    'BsmtFinSF1', 'BsmtFinSF2', 'BsmtUnfSF', 'TotalBsmtSF',
    'BsmtFullBath', 'BsmtHalfBath',
    'Fireplaces', 'PoolArea'
]
zero_impute_cols = [c for c in zero_impute_cols if c in X_train.columns]
print(f"5. Imputação de 0 para colunas numéricas com valores ausentes  ({len(zero_impute_cols)} colunas): {zero_impute_cols[:5]}...")

# 5. Colunas numéricas que receberão mediana
all_numerical = X_train.select_dtypes(include=['int64','float64']).columns.tolist()
median_impute_cols = [c for c in all_numerical if c not in zero_impute_cols and c != 'SalePrice']
print(f"6. Imputação da mediana para colunas numéricas ({len(median_impute_cols)} colunas): {median_impute_cols[:5]}...")

# 6. Colunas categóricas
categorical_cols = X_train.select_dtypes(include=['object','category']).columns.tolist()
print(f"7. Colunas categóricas ({len(categorical_cols)}): {categorical_cols[:5]}...")

# 7. Criar imputers
num_imputer_median = SimpleImputer(strategy='median')
num_imputer_zero   = SimpleImputer(strategy='constant', fill_value=0)
cat_imputer        = SimpleImputer(strategy='constant', fill_value='Missing')

# 8. ColumnTransformer
preprocessor = ColumnTransformer([
    ('median_impute', num_imputer_median, median_impute_cols),
    ('zero_impute',   num_imputer_zero,   zero_impute_cols),
    ('cat_impute',    cat_imputer,        categorical_cols)
], remainder='passthrough')

print("8. Pré-processamento das colunas pelos `SimpleImputer`s")
X_train_imputed = preprocessor.fit_transform(X_train)
X_test_imputed = preprocessor.transform(X_test)

print("9. Imputação terminada.")

# 10. Verificação
print(f"X_train dim. após imputação: {X_train_imputed.shape}")
print(f"X_test  dim. após imputação:  {X_test_imputed.shape}")
print("Término da primeira fase do pré-processamento.")
\end{verbatim}
\begin{enumerate}
\item Está Correto?
\label{sec:org2454961}
Sim, está correto. O log mostra exatamente o que se espera:

Dimensões:

Treino: (1022, 76) - pois de 81 colunas originais, removeu 4 (PoolQC, MiscFeature, Alley, Fence) e removeu o target SalePrice → 81 - 4 - 1 = 76.

Teste: (438, 76) - removeu as mesmas 4 colunas (teste não tem SalePrice) → 80 - 4 = 76.

Divisão das colunas:

Zero-impute: 12 colunas (ex.: GarageYrBlt, MasVnrArea).

Mediana: 24 colunas (ex.: LotFrontage).

Categóricas: 40 colunas (incluindo MSSubClass convertida para string).

Total: 12 + 24 + 40 = 76 → consistente.

Não há valores ausentes após a imputação (o código que você mostrou não exibe mais np.isnan, então provavelmente removemos essa verificação que dava erro de tipo).

Portanto, a primeira fase do pré-processamento (imputação) está perfeita.
\end{enumerate}
\subsection{Pré-Processamento dos Dados: \emph{Encoding} e \emph{Scaling}}
\label{sec:org50e751b}
Antes de produzir o modelo de regressão propriamente dito, ainda estamos interessados em duas estratégias para adequar os dados ao problema. Queremos \emph{codificar} (\texttt{encoding}) as colunas categóricas e \emph{escalar} (ou \emph{normalizar}, \textbf{scaling}) as colunas numéricas.

\begin{verbatim}
from sklearn.preprocessing import OneHotEncoder, StandardScaler
from sklearn.compose import ColumnTransformer
from sklearn.pipeline import Pipeline
from sklearn.linear_model import LinearRegression
from sklearn.metrics import mean_squared_error, mean_absolute_error, r2_score
import pandas as pd
import numpy as np

# 1. Verificação das listas de colunas e os DataFrames do bloco anterior
print("Colunas disponíveis :")
print(f"  median_impute_cols:   {len(median_impute_cols)}")
print(f"    zero_impute_cols:   {len(zero_impute_cols)}")
print(f"    categorical_cols:   {len(categorical_cols)}")
print(f"    total           :   {len(median_impute_cols) + len(zero_impute_cols) + len(categorical_cols)}")

# 2. Construir os pipelines para cada grupo e o transf. de colunas
numeric_median_pipe = Pipeline([
    ('impute_median', SimpleImputer(strategy='median')),
    ('scale', StandardScaler())
])

numeric_zero_pipe = Pipeline([
    ('impute_zero', SimpleImputer(strategy='constant', fill_value=0)),
    ('scale', StandardScaler())
])

categorical_pipe = Pipeline([
    ('impute_missing', SimpleImputer(strategy='constant', fill_value='Missing')),
    ('onehot', OneHotEncoder(handle_unknown='ignore', drop='first'))
])

# ColumnTransformer que aplica os três pipelines
preprocessor_final = ColumnTransformer([
    ('num_median', numeric_median_pipe, median_impute_cols),
    ('num_zero',   numeric_zero_pipe,   zero_impute_cols),
    ('cat',        categorical_pipe,    categorical_cols)
])

print("-" * 26);
print(preprocessor_final);

# 3. Transformar os dados
X_train_transformed = preprocessor_final.fit_transform(X_train, y_train)

print(X_train_transformed);
\end{verbatim}
\subsection{Definição do Modelo de Regressão Linear e Treinamento}
\label{sec:org5564d96}
A definição do modelo \emph{baseline} é bastante simples, tratando-se apenas de regressão linear conforme disponibilizado pelo \texttt{sklearn}:

\begin{verbatim}
model = LinearRegression()

# Usando log-transform no target
y_train_log = np.log1p(y_train)

model.fit(X_train_transformed, y_train_log)
\end{verbatim}
\subsection{Avaliação do Modelo \emph{Baseline}}
\label{sec:org46b34fe}

\begin{verbatim}
y_pred_log = model.predict(X_train_transformed)
y_pred = np.expm1(y_pred_log)

rmse = np.sqrt(mean_squared_error(y_train, y_pred))
mae  = mean_absolute_error(y_train, y_pred)
r2   = r2_score(y_train, y_pred)
rmse_log = np.sqrt(mean_squared_error(y_train_log, y_pred_log))

print("-" * 30)
print("11. Métricas no treino:")
print(f"RMSE (log): {rmse_log:.5f}")
print(f"RMSE: {rmse:.4f}")
print(f"MAE:  {mae:.4f}")
print(f"R²:   {r2:.4f}")

from sklearn.model_selection import cross_val_score

scores = -cross_val_score(
    model,
    X_train_transformed,
    y_train_log,
    scoring="neg_root_mean_squared_error",
    cv=5
)

print(f"CV RMSE (log): {scores.mean():.5f} ± {scores.std():.5f}")

y_test_pred_log = model.predict(X_test_transformed)
y_test_pred = np.expm1(y_test_pred_log)

print("12. Predições no conjunto de teste concluídas.")

submission = pd.DataFrame({
    "Id": test_df["Id"],
    "prediction": y_test_pred
})

submission.to_csv("submission.csv", index=False)
\end{verbatim}
\subsubsection{Resultados do Modelo \emph{Baseline}}
\label{sec:orgbff5ae6}
Métricas no treino:
\begin{itemize}
\item RMSE (log): 0.08926
\item RMSE: 17406.9043
\item MAE:  10889.9244
\item R²:   0.9497
\item CV RMSE (log): 0.18788 ± 0.05491
\end{itemize}

Em relação ao teste real, RMSE(log) \textasciitilde{}= 0.23
\section{Modelo \texttt{stacking-ensemble 4}}
\label{sec:org8c692e8}
\texttt{stacking-ensemble 4} é um \textbf{stacking ensemble} de dois níveis: cinco modelos de gradient boosting como camada base e um blend de RidgeCV + ElasticNet como meta-modelo.
\subsection{Pré-processamento}
\label{sec:org4aef768}
\subsubsection{Remoção de Outliers}
\label{sec:org87c76cf}
Antes de qualquer modelagem, quatro filtros são aplicados ao conjunto de treino para eliminar observações atípicas que distorceriam o aprendizado:

\begin{verbatim}
train_df = train_df[train_df["GrLivArea"] < 4500]
train_df = train_df[~((train_df["GrLivArea"] > 4000) & (train_df["SalePrice"] < 200000))]
train_df = train_df[train_df["LotArea"] < 100000]
train_df = train_df[train_df["TotalBsmtSF"] < 6000]
\end{verbatim}

O segundo filtro é especialmente importante: remove casas grandes \texttt{vendidas} a preço baixo, que são anomalias conhecidas do dataset Ames e prejudicam modelos de árvore.
\subsubsection{Transformação da Variável Alvo}
\label{sec:orge9fd220}

O target \texttt{SalePrice} é transformado com \texttt{log1p} antes do treino:

\begin{verbatim}
y = np.log1p(train_df["SalePrice"])
\end{verbatim}

Isso aproxima a distribuição da normal, reduz a influência de valores extremos e é o formato esperado na submissão — as predições são entregues em escala logarítmica, sem reverter com \texttt{expm1}.
\subsubsection{Imputação Contextual}
\label{sec:org9bd6a0c}
Em vez de imputação genérica por mediana, os valores ausentes são tratados com base no domínio do problema:

\begin{itemize}
\item \textbf{Categóricas de ausência semântica} (ex: \texttt{GarageType}, \texttt{PoolQC}): NaN
\end{itemize}
significa que o imóvel não possui aquela característica, portanto preenchidos com a string \texttt{"None"}.

\begin{itemize}
\item \textbf{Numéricas de área ou contagem} (ex: \texttt{GarageArea}, \texttt{BsmtFinSF1}):
\end{itemize}
imóveis sem garagem ou porão têm área zero, portanto preenchidos com \texttt{0}.

\begin{itemize}
\item \textbf{Demais numéricas}: imputação por mediana dentro de cada fold de validação cruzada, usando \texttt{SimpleImputer}.
\end{itemize}
\subsubsection{Feature Engineering}
\label{sec:orge54c152}
A função \texttt{add\_features} expande o conjunto original de 80 para 103 features, agrupadas em quatro categorias:
\begin{enumerate}
\item Agregações de Área
\label{sec:orgba0fdea}

\begin{verbatim}
df["TotalSF"]    = df["TotalBsmtSF"] + df["1stFlrSF"] + df["2ndFlrSF"]
df["TotalPorchSF"] = df["OpenPorchSF"] + df["EnclosedPorch"] + ...
df["TotalBathrooms"] = df["FullBath"] + 0.5*df["HalfBath"] + ...
\end{verbatim}

Capturam a área utilizável total, que é mais preditiva que qualquer área individual.
\item Flags Binárias de Presença
\label{sec:org577b4a4}

\begin{verbatim}
df["HasPool"]      = (df["PoolArea"] > 0).astype(int)
df["HasGarage"]    = (df["GarageArea"] > 0).astype(int)
df["HasBsmt"]      = (df["TotalBsmtSF"] > 0).astype(int)
df["HasFireplace"] = (df["Fireplaces"] > 0).astype(int)
df["IsRemodeled"]  = (df["YearRemodAdd"] != df["YearBuilt"]).astype(int)
df["IsNew"]        = (df["YrSold"] == df["YearBuilt"]).astype(int)
\end{verbatim}

Modelos de árvore beneficiam-se de splits binários explícitos quando a distinção presença/ausência é mais relevante que a magnitude.
\item Interações Multiplicativas
\label{sec:org0e2de67}

\begin{verbatim}
df["OverallQual_TotalSF"]   = df["OverallQual"] * df["TotalSF"]
df["OverallQual_GrLivArea"] = df["OverallQual"] * df["GrLivArea"]
df["OverallQual_sq"]        = df["OverallQual"] ** 2
df["QualCondScore"]         = df["OverallQual"] * df["OverallCond"]
\end{verbatim}

\texttt{OverallQual} é a variável mais correlacionada com preço. Suas
interações com área capturam o efeito não-linear de que qualidade
importa proporcionalmente mais em casas grandes.
\item Transformações Logarítmicas e Temporais
\label{sec:orge285d29}

\begin{verbatim}
df["LotAreaLog"]  = np.log1p(df["LotArea"])
df["GrLivAreaLog"]= np.log1p(df["GrLivArea"])
df["HouseAge"]    = df["YrSold"] - df["YearBuilt"]
df["RemodAge"]    = df["YrSold"] - df["YearRemodAdd"]
\end{verbatim}

\emph{Log-transforms} reduzem a \emph{skewness} de variáveis de área, ajudando especialmente o meta-modelo linear. As features de idade capturam depreciação e o efeito de reformas recentes.
\end{enumerate}
\subsection{Encoding de Variáveis Categóricas}
\label{sec:org49f4eb5}

As 44 variáveis categóricas são transformadas com \textbf{Target Encoding} dentro de cada fold:

\begin{verbatim}
enc = ce.TargetEncoder(cols=cat_cols, smoothing=10)
enc.fit(X_tr, y_tr)
\end{verbatim}

O parâmetro \texttt{smoothing=10} é crítico: mistura a média da categoria com a média global proporcionalmente ao tamanho da categoria, evitando \emph{data leakage} em categorias raras. O valor 10 (versus 0.3 do baseline) reduz significativamente o \emph{overfitting} do \emph{encoding}.

O encoder é sempre \emph{fitado} apenas nos dados de treino do fold e aplicado em validação e teste, respeitando o isolamento de
informação da validação cruzada.
\subsection{Definição de Modelos Base}
\label{sec:orgea9ca8f}
Cinco modelos de \emph{gradient boosting} com hiperparâmetros distintos formam a camada base do \emph{ensemble}. A diversidade entre eles é intencional para que erros não correlacionados se cancelam no stacking.

\begin{center}
\begin{tabular}{lrrrl}
Modelo & \texttt{n\_estimators} & \texttt{learning\_rate} & \texttt{num\_leaves=/=depth} & Característica\\
\hline
LGB1 & 6000 & 0.005 & 31 & Conservador, balanceado\\
LGB2 & 8000 & 0.003 & 15 & Raso, muito regularizado\\
LGB3 & 5000 & 0.008 & 63 & Muitas folhas, agressivo\\
Cat & 5000 & 0.007 & 6 & Melhor OOF individual\\
XGB & 5000 & 0.005 & 5 & Complementa LGB/Cat\\
\end{tabular}
\end{center}

OOF RMSE individuais (seed 42):
\begin{itemize}
\item LGB1: 0.126874
\item LGB2: 0.123090
\item LGB3: 0.119819
\item CatBoost: \textbf{0.116575} (melhor)
\item XGBoost: 0.125426
\end{itemize}
\subsection{Validação Cruzada (Out-of-Fold)}
\label{sec:orgdf7697c}

\begin{verbatim}
kf = KFold(n_splits=10, shuffle=True, random_state=42)
\end{verbatim}

10 folds foram utilizados em vez dos 5 do baseline por dois motivos: cada modelo é avaliado em 10 conjuntos distintos, reduzindo a variância da estimativa de RMSE; e cada modelo de treino usa 90\% dos dados em vez de 80\%, gerando predições mais fortes.

As predições \emph{Out-of-Fold} (OOF) são usadas como features para o meta-modelo, garantindo que o meta-modelo nunca veja predições feitas sobre dados em que o modelo base foi treinado.
\subsection{Meta-Modelo (\emph{Stacking})}
\label{sec:org9e1f63c}
As predições OOF dos cinco modelos base formam uma matriz de entrada para dois meta-modelos:

\begin{verbatim}
# Meta 1
meta_ridge = RidgeCV(alphas=np.logspace(-4, 2, 50), cv=10)
meta_ridge.fit(X_meta, y)   # RMSE: 0.113588

# Meta 2
meta_en = ElasticNet(alpha=0.0003, l1_ratio=0.5, max_iter=50000)
meta_en.fit(X_meta, y)      # RMSE: 0.114635

# Blend final (50/50)
final = 0.5 * meta_ridge.predict(...) + 0.5 * meta_en.predict(...)
# RMSE: 0.113965
\end{verbatim}

\texttt{RidgeCV} seleciona automaticamente o melhor \texttt{alpha} entre 50 valores via \emph{cross-validation} interna. O \emph{blend} 50/50 com \texttt{ElasticNet} reduz ligeiramente a variância da predição final em relação a usar apenas um dos dois.
\subsection{Resultados para o Modelo \texttt{stacking-ensemble 4}}
\label{sec:orgd89f27f}
\begin{itemize}
\item OOF RMSE FINAL (blend meta):  0.113965
\item OOF RMSE FINAL (média direta): 0.119427
\end{itemize}
\subsubsection{Decisões de Design para o Modelo \texttt{stacking-ensemble 4}}
\label{sec:org109fed3}
\begin{center}
\begin{tabular}{lll}
Decisão & Alternativa descartada & Motivo\\
\hline
Imputação contextual por domínio & Mediana global & NaN em GarageArea significa 0, não mediana\\
RobustScaler & StandardScaler & Mais estável com outliers remanescentes\\
TargetEncoding smoothing=10 & smoothing=0.3 & Menos leakage em categorias raras\\
10-fold CV & 5-fold & Menor variância, mais dados por treino\\
5 modelos heterogêneos & Apenas LightGBM & Erros não correlacionados melhoram o blend\\
RidgeCV como meta & ElasticNet puro & Seleção automática de alpha, mais estável\\
Predição em log scale & \texttt{expm1} antes de submeter & Formato exigido pela competição\\
\end{tabular}
\end{center}
\end{document}
